\begin{helper}
Assume \(\noderef{ParameterHierarchy}(\varepsilon,\varepsilon_1,\varepsilon_2,n_0)\).
For every \(\delta>0\) there are a positive internal error \(\eta>0\) and
a threshold \(N\ge n_0\) such that, for every \(n>N\), if
\[
K=\noderef{TwoBiteNaturalIndependenceScale}(1+\varepsilon,n),
\qquad
p=\noderef{TwoBiteEdgeProbability}\!\left(\frac12,n\right),
\]
then
\[
\sum_{\ell_R=0}^{K}\sum_{\ell_B=0}^{K}
{\bf 1}_{\ell_R\ell_B\ge K}\,
\exp\!\left(
\left(\eta-\frac{2-\ell_R/K-\ell_B/K}{2}\right)K\log n
\right)
\min\!\left\{1,
\exp\!\left(-p\bigl(f(\ell_R,\ell_B)-\varepsilon^3K^2\bigr)\right)
\right\}
\le
\delta\binom nK^{-1},
\]
where \(f=\noderef{ProjectionOpenPairFunction}\), all cardinalities are cast to
real numbers in the displayed expression, and the summand is \(0\) when
\(\ell_R\ell_B<K\).
\end{helper}
\begin{proof}
From \(\noderef{ParameterHierarchy}\), we have
\[
0<\varepsilon<\frac1{16}.
\]
Choose \(\eta>0\) small enough that
\[
\eta\le \frac{\varepsilon}{80},\qquad
\eta\le \frac{\varepsilon^2}{320},\qquad
\eta\le
\frac{\varepsilon(1-\varepsilon-22\varepsilon^2)}
     {320(1+\varepsilon)}.
\tag{1}
\]
The last upper bound is positive because
\(\varepsilon+22\varepsilon^2<1\) for \(\varepsilon<1/16\).

Apply \(\noderef{FixedSetProjectionBinomialSummationMargins}\) with this
\(\eta\) and the given \(\delta\), and apply the upper- and middle-regime
children
\(\noderef{FixedSetProjectionUpperRegimeExponent}\) and
\(\noderef{FixedSetProjectionMiddleRegimeExponent}\) with the same \(\eta\).
Let \(N\ge n_0\) be larger than the three thresholds supplied by these children.
For \(n>N\), put
\[
K=\noderef{TwoBiteNaturalIndependenceScale}(1+\varepsilon,n),
\qquad
p=\noderef{TwoBiteEdgeProbability}\!\left(\frac12,n\right).
\]
The binomial-summation child gives
\[
1\le K\le n,\qquad
\log\binom nK\le \left(\frac12+\eta\right)K\log n, \tag{2}
\]
and
\[
2\log(K+1)+\log(\delta^{-1})\le \eta K\log n. \tag{3}
\]
Since \(1\le K\le n\), increasing \(N\) if necessary gives
\[
0\le K\log n. \tag{4}
\]

Fix a nonzero summand, so \(\ell_R\ell_B\ge K\), and put
\[
x_R=\ell_R/K,\qquad x_B=\ell_B/K.
\]
Since \(1\le K\) and \(0\le\ell_R,\ell_B\le K\), we have
\[
0\le x_R,x_B\le1.
\]
The summand multiplied by \(\binom nK\) has logarithm at most
\[
\left(\eta+\left(\frac12+\eta\right)-\frac{2-x_R-x_B}{2}\right)K\log n
\]
if the minimum contributes \(1\).  If the conditional exponential contributes,
the same expression has the additional term
\[
-p\bigl(f(\ell_R,\ell_B)-\varepsilon^3K^2\bigr).
\]

First suppose \(x_R+x_B\le1-\varepsilon/2\).  The minimum is at most \(1\), and
\(\noderef{FixedSetProjectionLowerRegimeExponent}\) applies with the natural
boundary data \(1\le K\) from (2), \(0\le\ell_R,\ell_B\le K\) from the summation
range, \(0\le K\log n\) from (4), and \(\eta\le\varepsilon/80\) from (1).  In
the logarithmic exponent of \(\binom nK\) times the summand, the minimum factor
contributes at most
\[
\min\left\{0,\,
-p\left(f(\ell_R,\ell_B)-\varepsilon^3K^2\right)\right\}.
\]
The lower-regime child therefore gives
\[
\log\!\left(\binom nK\cdot\hbox{summand}\right)\le -4\eta K\log n. \tag{5}
\]

Next suppose \(x_R+x_B\ge1+\varepsilon/2\).  In the definition of
\(f=\noderef{ProjectionOpenPairFunction}\), both minima are at most the first
listed term, so
\[
f(\ell_R,\ell_B)
\ge
\binom{\ell_R}{2}_{\mathbb R}+\binom{\ell_B}{2}_{\mathbb R}
-\binom{K-\ell_R}{2}_{\mathbb R}
-\binom{K-\ell_B}{2}_{\mathbb R}.
\]
The upper-regime child then gives
\[
\log\!\left(\binom nK\cdot\hbox{summand}\right)
\le
-\frac{\varepsilon(x_R+x_B-1-\varepsilon^2-\varepsilon^3)}{2}K\log n
+3\eta K\log n.
\]
Because \(x_R+x_B\ge1+\varepsilon/2\) and \(\varepsilon<1/16\),
\[
-\frac{\varepsilon(x_R+x_B-1-\varepsilon^2-\varepsilon^3)}{2}
\le -\frac{\varepsilon^2}{8}.
\]
Using \(\eta\le\varepsilon^2/320\), we have
\[
-\frac{\varepsilon^2}{8}+3\eta\le -4\eta,
\]
because \(7\eta\le\varepsilon^2/8\).  With (4), the upper-regime exponent is
therefore at most
\[
-4\eta K\log n. \tag{6}
\]

It remains to treat
\[
1-\varepsilon/2<x_R+x_B<1+\varepsilon/2.
\]
If \(x_B\le x_R\), then the middle-regime child gives \(pn<K-\ell_B\) and
\[
\begin{aligned}
&\left(\eta+\left(\frac12+\eta\right)-\frac{2-x_R-x_B}{2}\right)K\log n\\
&\quad
-p\left(
\binom{\ell_R}{2}_{\mathbb R}+\binom{\ell_B}{2}_{\mathbb R}
-\binom{K-\ell_R}{2}_{\mathbb R}
-\binom{pn}{2}_{\mathbb R}
-\binom{K-\ell_B-pn}{2}_{\mathbb R}
-\varepsilon^3K^2\right)\\
&\le
\frac{\varepsilon(-1+\varepsilon+22\varepsilon^2)}{16(1+\varepsilon)}
K\log n+3\eta K\log n.
\end{aligned}
\]
The displayed upper bound applies to the exponent of
\(\binom nK\) times the summand because the second minimum in \(f\) is bounded
above by
\[
\binom{pn}{2}_{\mathbb R}+\binom{K-\ell_B-pn}{2}_{\mathbb R}
\]
and the first by \(\binom{K-\ell_R}{2}_{\mathbb R}\).  If \(x_R\le x_B\), the
same argument applies after interchanging red and blue.  Since
\[
\frac{\varepsilon(-1+\varepsilon+22\varepsilon^2)}{16(1+\varepsilon)}
=-\frac{\varepsilon(1-\varepsilon-22\varepsilon^2)}{16(1+\varepsilon)}
\]
and \(\eta\) satisfies (1), the right side above plus \(3\eta\) is at most
\(-4\eta\): indeed the bound on \(\eta\) gives
\[
\frac{\varepsilon(1-\varepsilon-22\varepsilon^2)}{16(1+\varepsilon)}
\ge 20\eta.
\]
With (4), this middle-regime exponent is also at most
\[
-4\eta K\log n. \tag{7}
\]

The bounds (5), (6), and (7) show that every nonzero summand is at most
\[
\binom nK^{-1}\exp(-4\eta K\log n).
\]
There are at most \((K+1)^2\) summands, and (3) together with (4) gives
\[
(K+1)^2\exp(-4\eta K\log n)
\le
\delta\exp(-3\eta K\log n)
\le \delta.
\]
Multiplying by \(\binom nK^{-1}\) proves the claimed finite sum bound.  The
summands with \(\ell_R\ell_B<K\) are \(0\) by definition, so they do not affect
the estimate.
\end{proof}
